\documentclass[11pt,a4paper]{article}
\usepackage[margin=2.5cm]{geometry}
\usepackage{amsmath,amssymb}
\usepackage{enumitem}
\usepackage[hidelinks]{hyperref}

\newcommand{\given}{\,|\,}
\newcommand{\bin}{\operatorname{binomial}}
\newcommand{\pois}{\operatorname{Poisson}}
\newcommand{\Emax}{E_{\max}}
\newcommand{\Imax}{I_{\max}}

\title{Allowing sporadic downward revisions of reported case counts\\
\large An extension of the renewal-equation reporting-delay model}
\author{}
\date{\today}

\begin{document}
\maketitle

\section{Motivation}

The model of Bajaj, Thompson and Lambert \cite{bajaj2025} assumes that the
reported count of cases with symptom onset on day $t$, when viewed from a later
day $t'$, satisfies $C(t\given t')\le I(t)$ and can only increase as $t'$
increases. Both properties follow from equation (2.2) of that paper,
\begin{equation}\label{eq:original}
C(t\given t'') - C(t\given t') \sim
\bin\bigl(I(t) - C(t\given t'),\; p_\theta(t, t''-t'\given t')\bigr),
\qquad t''>t'.
\end{equation}
In practice, counts are occasionally revised downwards, for example when a
false positive or a duplicate is removed or when a case is reassigned to a
different onset day. Here we extend the model so that such revisions have
non-zero probability, while recovering the original model as a special case.
Notation follows \cite{bajaj2025}; new symbols are introduced explicitly.

\section{Notation}

\begin{itemize}[itemsep=2pt]
\item $t=1,\dots,T$ indexes onset days.
\item $t_0<t_1<\dots<t_J$ are the times at which the count for onset day $t$
is observed. All satisfy $t_j\ge t$, and they may differ between onset days.
\item $C(t\given t_j)$ is the reported count for onset day $t$ as seen at time
$t_j$ (data).
\item $I(t)$ is the true count of cases with onset on day $t$ (latent).
\item $S_\theta(t\given t')$ is the survival function of the reporting delay
distribution with parameters $\theta$, and
$p_\theta(t,t''-t'\given t')=\bigl[S_\theta(t\given t')-S_\theta(t\given t'')\bigr]/S_\theta(t\given t')$.
\end{itemize}
As in \cite{bajaj2025}, the first observation $C(t\given t_0)$ is conditioned
on, and the likelihood is built from changes between successive observations.

\section{Model}

\subsection{Latent decomposition}

For each onset day $t$ and observation time $t_j$ we write
\begin{equation}\label{eq:decomp}
C(t\given t_j) = D(t\given t_j) + E(t\given t_j),
\end{equation}
where $D(t\given t_j)$ is the number of records held at $t_j$ for onset day $t$
that are genuine cases with onset on day $t$, and $E(t\given t_j)$ is the number
of records held at $t_j$ for onset day $t$ that are erroneous and have not yet
been withdrawn. Neither is observed.

\subsection{Increments}

Two new parameters are introduced: $\eta\ge 0$, the expected number of
erroneous records entering per genuine report, and $\rho\in(0,1)$, the daily
probability that an erroneous record is withdrawn. For onset day $t$ and the
interval $(t_{j-1},t_j]$, $j=1,\dots,J$, define
\begin{align}
\Delta C_j(t) &:= C(t\given t_j) - C(t\given t_{j-1}), \\
p_j(t) &:= p_\theta(t,\, t_j-t_{j-1}\given t_{j-1}), \\
\lambda_j(t) &:= \eta\, I(t)\,\bigl[S_\theta(t\given t_{j-1}) - S_\theta(t\given t_j)\bigr], \\
r_j &:= 1-(1-\rho)^{\,t_j-t_{j-1}}.
\end{align}
Note that $\Delta C_j(t)$ may now be negative. Within the interval, three
independent increments occur, all pertaining to onset day $t$:
\begin{align}
d_j(t) &\sim \bin\bigl(I(t) - D(t\given t_{j-1}),\; p_j(t)\bigr)
  && \text{genuine cases newly reported,} \label{eq:d}\\
n_j(t) &\sim \pois\bigl(\lambda_j(t)\bigr)
  && \text{erroneous records newly entered,} \label{eq:n}\\
m_j(t) &\sim \bin\bigl(E(t\given t_{j-1}),\; r_j\bigr)
  && \text{erroneous records withdrawn,} \label{eq:m}
\end{align}
so that
\begin{align}
D(t\given t_j) &= D(t\given t_{j-1}) + d_j(t), \\
E(t\given t_j) &= E(t\given t_{j-1}) + n_j(t) - m_j(t), \\
\Delta C_j(t) &= d_j(t) + n_j(t) - m_j(t). \label{eq:deltaC}
\end{align}

\paragraph{Conventions.}
\begin{enumerate}[label=(\roman*),itemsep=2pt]
\item Only erroneous records already present at $t_{j-1}$ can be withdrawn in
$(t_{j-1},t_j]$; a record cannot enter and be withdrawn within the same
interval.
\item Conditional on $\{I(t)\}_{t=1}^T$, $\theta$, $\eta$ and $\rho$,
increments are independent across onset days.
\item $E(t\given t_0)=0$, so that $D(t\given t_0)=C(t\given t_0)$. This is an
approximation: it assumes that no erroneous records are present at the first
observation.
\end{enumerate}

\paragraph{Properties.}
Equation \eqref{eq:d} is equation \eqref{eq:original} applied to $D$ rather
than $C$, so $D(t\given t_j)\le I(t)$ and $D$ is non-decreasing in $t_j$,
whereas $C(t\given t_j)$ may exceed $I(t)$ and may decrease. Setting $\eta=0$
gives $E\equiv 0$ and $C\equiv D$, which recovers the original model. Because
erroneous records are eventually withdrawn, $E(t\given t')\to 0$ as
$t'\to\infty$, and the identity $I(t)=C(t\given\infty)$ is retained. The mean
$\lambda_j(t)$ is $\eta$ times the expected number of genuine cases with onset
on day $t$ that are reported in $(t_{j-1},t_j]$, so erroneous records
concentrate shortly after onset and tail off as the delay distribution does.

\section{Likelihood}

Since $E(t\given t_j)$ is unobserved, we sum over it. Downward revisions are
assumed sporadic, so $E$ is truncated at a small value $\Emax$. For
$e=0,\dots,\Emax$ define the forward probabilities
\begin{equation}
\alpha_j(e;t) := P\bigl(C(t\given t_1),\dots,C(t\given t_j),\; E(t\given t_j)=e
\;\big|\; I(t),\, C(t\given t_0),\, \theta,\eta,\rho\bigr).
\end{equation}
These satisfy $\alpha_0(e;t)=\mathbf{1}\{e=0\}$ and, for $j=1,\dots,J$,
\begin{equation}\label{eq:forward}
\begin{aligned}
\alpha_j(e;t) = \sum_{e'=0}^{\Emax} \alpha_{j-1}(e';t) \sum_{m=0}^{e'}
\;&\bin\bigl(m \given e',\, r_j\bigr) \\
\times\;&\pois\bigl(e-e'+m \given \lambda_j(t)\bigr) \\
\times\;&\bin\bigl(\Delta C_j(t)-(e-e') \given I(t)-C(t\given t_{j-1})+e',\; p_j(t)\bigr),
\end{aligned}
\end{equation}
where $\operatorname{dist}(x\given\phi)$ denotes the probability mass function
of $\operatorname{dist}$ evaluated at $x$. In the summand, with
$E(t\given t_{j-1})=e'$ and $E(t\given t_j)=e$: $m$ is $m_j(t)$; $e-e'+m$ is
$n_j(t)$; $\Delta C_j(t)-(e-e')$ is $d_j(t)$, by \eqref{eq:deltaC}; and
$I(t)-C(t\given t_{j-1})+e'$ is $I(t)-D(t\given t_{j-1})$, by \eqref{eq:decomp}.
Any term involving a negative count, or a binomial count exceeding its size,
has probability zero. The likelihood contribution of onset day $t$ is
\begin{equation}\label{eq:lik}
L_t\bigl(I(t);\theta,\eta,\rho\bigr) := \sum_{e=0}^{\Emax}\alpha_J(e;t),
\end{equation}
at a cost of $O(J\,\Emax^3)$ per value of $I(t)$.

\section{Changes to the sampling algorithm}

We describe the changes to the Poisson model of \cite{bajaj2025}; the
unrestricted model is modified analogously.

\paragraph{Sampling $R_t$.} Equation (2.4) of \cite{bajaj2025} is unchanged,
since it depends on the data only through $\{I(t)\}$.

\paragraph{Sampling the true case counts.} The product of binomial terms in
equation (2.6) of \cite{bajaj2025} is replaced by \eqref{eq:lik}. For
$I(t)\in\{0,1,\dots,\Imax\}$,
\begin{equation}
\begin{aligned}
q\bigl(I(t)\bigr) ={}& L_t\bigl(I(t);\theta,\eta,\rho\bigr)
\times \pois\Bigl(I(t)\,\Big|\, R_{i[t]}\sum_{s=1}^{t-1}\omega_s I(t-s)\Bigr) \\
&\times \prod_{\tau=t+1}^{T}\pois\Bigl(I(\tau)\,\Big|\, R_{i[\tau]}\sum_{s=1}^{\tau-1}\omega_s I(\tau-s)\Bigr).
\end{aligned}
\end{equation}
The lower bound $I(t)\ge C(t\given t_J)$ in equation (2.5) of \cite{bajaj2025}
is removed; values of $I(t)$ that are too small to be consistent with the data
receive zero probability through \eqref{eq:forward}. The update remains a draw
from a one-dimensional discrete distribution.

\paragraph{Sampling the reporting parameters.} Equation (2.7) of
\cite{bajaj2025} becomes
\begin{equation}
p\bigl(\theta,\eta,\rho \,\big|\, \{I(t)\}_{t=1}^{T},\{C(t\given t_j)\}\bigr)
\propto p(\theta)\,p(\eta)\,p(\rho)\prod_{t=1}^{T}L_t\bigl(I(t);\theta,\eta,\rho\bigr),
\end{equation}
which is sampled by random-walk Metropolis as before, with $\eta$ and $\rho$
added to the block.

\section{Remarks}

\begin{itemize}[itemsep=2pt]
\item \emph{Identifiability.} The parameters $\eta$ and $\rho$ are informed
only by the rare intervals in which $\Delta C_j(t)<0$. An erroneous record not
withdrawn by $t_J$ is indistinguishable from a genuine case. Informative priors
concentrated near zero are therefore advisable (for example
$\eta\sim\operatorname{beta}(1,500)$), and it may be preferable to fix $\rho$.
\item \emph{Nowcasts.} For recent onset days some of $C(t\given t_J)$ may be
erroneous, so the posterior for $I(t)$ is shifted slightly downwards and
widened relative to the original model.
\item \emph{Reassigned onset days.} A case moved from onset day $t$ to onset
day $t^*$ is represented marginally, as an erroneous record on day $t$ and a
late genuine report on day $t^*$. Modelling the two jointly would couple onset
days and require block updates of $\bigl(I(t),I(t^*)\bigr)$.
\item \emph{Truncation.} $\Emax$ should be chosen so that the posterior
probability of $E(t\given t_j)=\Emax$ is negligible; this can be checked from
the forward probabilities.
\end{itemize}

\begin{thebibliography}{1}
\bibitem{bajaj2025}
Bajaj S, Thompson R, Lambert B. 2025
A renewal-equation approach to estimating $R_t$ and infectious disease case
counts in the presence of reporting delays.
\emph{Phil. Trans. R. Soc. A} \textbf{383}, 20240357.
\href{https://doi.org/10.1098/rsta.2024.0357}{doi:10.1098/rsta.2024.0357}
\end{thebibliography}

\end{document}
